\documentclass[12pt]{ltjsarticle}

\usepackage{luatexja}
\usepackage{amsmath,amssymb}
\usepackage{longtable,booktabs,array}

\title{sample}
\author{隼人 今村}
\date{April 2026}

\begin{document}

\maketitle

\section{班員の紹介}

{\def\LTcaptype{none}
\begin{longtable}[]{@{}lllll@{}}
\toprule\noalign{}
\endhead
\bottomrule\noalign{}
\endlastfoot
名前 & 学年 & 出身 & 気になっていること & つくば市のおすすめスポット \\
小倉 & 3 & 神奈川県 & パズル & 筑波山 \\
櫻井 & 3 & 茨城県 & ポーカー & えびすや \\
村田 & 3 & 埼玉県 & ロードバイク & りんりんロード \\
\end{longtable}
}

\section{階層分析}

階層分析法（AHP）は、複数の基準をもつ意思決定問題を
\textbf{階層構造}・\textbf{一対比較}・\textbf{総合評価}
によって整理する手法である。

\subsection{基本の流れ}

\begin{itemize}
\item 目的を定める
\item 評価基準を設定する
\item 代替案を比較する
\item 重みを求めて総合評価する
\end{itemize}

\subsection{階層構造の例}

\[
\text{目的} \rightarrow \text{評価基準} \rightarrow \text{代替案}
\]

\begin{itemize}
\item 第1階層：目的
\item 第2階層：評価基準
\item 第3階層：代替案
\end{itemize}

\subsection{一対比較}

比較行列を \(A = \left( a_{ij} \right)\) とすると、各成分は対象 \(i\)
の対象 \(j\) に対する重要度を表す。

\[
A = \begin{pmatrix}
1 & a_{12} & \ldots & a_{1n} \\
\frac{1}{a_{12}} & 1 & \ldots & a_{2n} \\
\vdots & \vdots & \ddots & \vdots \\
\frac{1}{a_{1n}} & \frac{1}{a_{2n}} & \ldots & 1
\end{pmatrix}
\]

性質は次の通りである。

\begin{itemize}
\item \(a_{ii} = 1\)
\item \(a_{ji} = \frac{1}{a_{ij}}\)
\item 比較尺度には \(1,3,5,7,9\) とその逆数を用いる
\end{itemize}

\subsection{重みの計算}

\subsubsection{固有値法}

最大固有値に対応する固有ベクトルを用いて重みを求める。

\[
Aw = \lambda_{\text{max}} w
\]

\[
\sum_{i=1}^{n} w_i = 1
\]

\subsubsection{幾何平均法}

各行の幾何平均を使って重みを求める。

\[
g_i = \sqrt[n]{\prod_{j=1}^{n} a_{ij}}
\]

\[
w_i = \frac{g_i}{\sum_{j=1}^{n} g_j}
\]

\subsection{整合性指標}

比較の一貫性は CI によって確認する。

\[
CI = \frac{\lambda_{\text{max}} - n}{n - 1}
\]

\begin{itemize}
\item \(CI = 0\) に近いほど整合性が高い
\item 大きい場合は再評価が必要になる
\end{itemize}

\subsection{例：今日の献立}

評価基準と代替案を次のようにする。

{\def\LTcaptype{none}
\begin{longtable}{@{}p{0.38\linewidth}p{0.28\linewidth}p{0.28\linewidth}@{}}
\toprule\noalign{}
項目 & 内容 & 例 \\
\midrule\noalign{}
\endhead
\bottomrule\noalign{}
\endlastfoot
目的 & 献立を選ぶ & 今日の献立 \\
基準 & 値段・満足感・季節感 & 3基準 \\
代替案 & 候補メニュー & カレー / カツ丼 / ハンバーグ \\
\end{longtable}
}

基準の重みを次のようにする。

{\def\LTcaptype{none}
\begin{longtable}[]{@{}llll@{}}
\toprule\noalign{}
値段 & ガッツリ感 & 季節感 & 合計 \\
\midrule\noalign{}
\endhead
\bottomrule\noalign{}
\endlastfoot
\(0.19\) & \(0.08\) & \(0.73\) & \(1.00\) \\
\end{longtable}
}

各代替案の総合評価を \(W_t\) とすると、重み付き和で表せる。

\[
W_t = W_p w_t^p + W_q w_t^q + W_r w_t^r
\]

評価例を表にすると次のようになる。

{\def\LTcaptype{none}
\begin{longtable}[]{@{}lllll@{}}
\toprule\noalign{}
代替案 & 値段 & ガッツリ感 & 季節感 & 総合評価 \\
\midrule\noalign{}
\endhead
\bottomrule\noalign{}
\endlastfoot
カレーライス & \(0.73\) & \(0.26\) & \(0.66\) & \(0.64\) \\
カツ丼 & \(0.19\) & \(0.64\) & \(0.16\) & \(0.20\) \\
ハンバーグ & \(0.08\) & \(0.10\) & \(0.18\) & \(0.16\) \\
\end{longtable}
}

\subsection{まとめ}

AHP では、次のような流れで意思決定を整理できる。

\begin{enumerate}
\item 階層を作る
\item 一対比較を行う
\item 重みを求める
\item 総合評価で順位を決める
\end{enumerate}
\subsection{フォントの変更}

フォントの変更（{\rmfamily Roman}）

フォントの変更（{\sffamily Sans Serif}）

フォントの変更（{\ttfamily Typewriter}）

\subsection{フォントサイズの変更}

フォントサイズの変更（{\small 小さい文字}）

フォントサイズの変更（{\Large 大きい文字}）

フォントサイズの変更（{\Huge とても大きい文字}）

\end{document}